\documentclass[10pt,a4paper]{amsart}
\usepackage[margin=19mm]{geometry}
\usepackage{amsmath,amssymb,mathtools}
\usepackage[scr=rsfs,cal=euler]{mathalfa}
\usepackage{xfrac}
\usepackage[unicode,hidelinks,bookmarks=false]{hyperref}
\newcommand{\ie}{i.e.\ }
\newcommand{\cf}{cf.\ }
\newcommand{\resp}{respectively\ }
\newcommand{\Subsection}{\subsection}
\providecommand{\nobreakdash}{\nobreak\hbox{-}}
\setlength{\emergencystretch}{2em}
\multlinegap0pt
\usepackage{bm}
\usepackage{bbm}
\usepackage{dsfont}
\usepackage{xspace}
\usepackage{cancel}
\usepackage{mathtools}
\usepackage{amsthm}
\makeatletter
\@ifundefined{assumption}{\newtheorem{assumption}{Assumption}}{}
\@ifundefined{lemma}{\newtheorem{lemma}{Lemma}}{}
\makeatother
\makeatletter
\newcommand{\E}{\mathbb{E}}
\newcommand{\R}{\mathbb{R}}
\newcommand{\be}{\begin{equation}}
\newcommand{\ee}{\end{equation}}
\expandafter\ifx\csname proof\endcsname\relax
\newenvironment{proof}{{\noindent\it Proof.}\quad}{\hfill $\square$\par}
\fi
\makeatother
\title[Adaptive directional decomposition methods for nonconvex con]{Adaptive directional decomposition methods for nonconvex constrained optimization}
\author{Qiankun Shi, Xiao Wang}
\date{Source: arXiv:2511.03210v2 (CC BY 4.0)}
\begin{document}
\maketitle

\noindent\emph{Source and licence.} Adaptive directional decomposition methods for nonconvex constrained optimization. Version arXiv:2511.03210v2 (CC BY 4.0), distributed under the Creative Commons Attribution 4.0 International licence, as recorded in the arXiv licence metadata for this exact version and shown on the abstract page https://arxiv.org/abs/2511.03210v2. Full rights notice: licenses/arxiv-2511.03210-CC-BY-4.0.txt.\par\smallskip

\noindent\textit{Editorial scope.} Counted unit: the Lemma whose source label is \texttt{lem:gradient\_estimator\_recursive\_lemma} in the article (source0.tex lines 1606--1623, source label \texttt{lem:gradient\_estimator\_recursive\_lemma}), delivered together with the article's own complete original proof of it (source0.tex lines 1625--1679, one \texttt{proof} environment of 55 lines). No mathematics is written, repaired or re-derived: statement and proof are the article's own text line for line. Required context retained uncounted: ass:stochastic (lines 1482--1484); ass:ms (lines 1590--1592); lm:vazuma (lines 2195--2200). Every definition, display and label of those lines is retained unchanged. Omitted from this package and not redistributed: 1--1481, 1485--1589, 1593--1605, 1624--1624, 1680--2194, 2201--2585. Nothing inside a retained excerpt is omitted or altered: every retained body is delivered exactly as the article's lines are written, so no omission receipt of any kind accompanies this package. Citations: the retained excerpts carry no citation command at all, so no bibliography entry of the article is transcribed and no reference is redistributed. Reference adaptation: none was needed and none was made; every reference inside the delivered unit resolves inside it, so no unresolved reference or citation is produced. Printed slips kept literal: no printed word, symbol, hypothesis or number of the counted statement or of its proof was altered. Layout and class: the article's own class is \texttt{article} with packages including amsmath, graphicx, listings, caption, subfigure, lineno, bm, subeqnarray, amsfonts, amssymb, amsbsy, color, lipsum, booktabs; the wrapper is the lane's pinned \texttt{amsart} editorial wrapper at 10pt on A4, which restates the theorem-like environments the retained text needs and adds the article's own macro definitions that the retained text needs (\texttt{\textbackslash E}, \texttt{\textbackslash R}, \texttt{\textbackslash be}, \texttt{\textbackslash ee}), with extra wrapper packages (bm, bbm, dsfont, xspace, cancel, mathtools). The delivered numbering is the unit's own. Assets: no retained span contains a figure environment, a graphics-inclusion command, a PDF-page inclusion, a table or a TikZ picture; no image file is copied into this package, the package declares no asset dependency and no image file is redistributed with it. Licence: version arXiv:2511.03210v2 of this article is distributed under the Creative Commons Attribution 4.0 International licence, as recorded in the arXiv licence metadata for this exact version and shown on the abstract page https://arxiv.org/abs/2511.03210v2. A full rights notice is delivered at licenses/arxiv-2511.03210-CC-BY-4.0.txt. Characters: every retained excerpt is pure ASCII, so the delivered engine, XeLaTeX with its default Latin Modern text font, reports no missing character.
\begin{assumption}\label{ass:stochastic}
{For any $x \in \R^d$, we have $\E_\xi[\nabla F(x,\xi)] = \nabla f(x)$, $\E_\zeta[C(x,\zeta)] = c(x)$ and $\E_\zeta[\nabla C(x,\zeta)] = \nabla c(x)$, and for almost {any} $\xi$ and $\zeta$, \(\|\nabla F(x,\xi) - \nabla f(x)\| \leq \sigma_f\), \(\|C(x,\zeta) - c(x)\| \leq \sigma_v\) and \(\|\nabla C_i(x,\zeta) - \nabla c_i(x)\| \leq \sigma_c\), $i=1,\ldots,m$.} 
\end{assumption}

\begin{assumption}\label{ass:ms}
For almost any $\xi\in \Xi$, $F(\cdot,\xi)$ has $L_g^f$-Lipschitz continuous gradients. For almost any  $\zeta\in \Xi$, $C_i(\cdot,\zeta)$, \(i = 1,\ldots,m\) are $L_c $-Lipschitz continuous and have  $L_g^c$-Lipschitz continuous gradients.
\end{assumption}

\begin{lemma}[Vector Azuma-Hoeffding inequality]\label{lm:vazuma}
Let $\{u_k\}$ be a vector-valued martingale difference, where $\E[u_k|\mathcal F(u_1,...,u_{k-1})] = 0$ and $\|u_k\| \leq a_k$. Then with probability at least $1-\gamma$ it holds that
\begin{align*}
    \bigg\|\sum_k u_k\bigg\| \leq 3\sqrt{\log(1/\gamma)\sum_k a_k^2}.
\end{align*}
\end{lemma}

\begin{lemma}\label{lem:gradient_estimator_recursive_lemma}
Under Assumptions \ref{ass:stochastic} and \ref{ass:ms}, given \(\gamma \in (0,1)\), suppose that $\{B_k\}$ satisfies 
\begin{align}\label{rm_b0}
    B_0^f =  \frac{81\sigma_f^2\log^2(1 /\gamma)}{\bar \epsilon_f^2},~ B_0^c =  \frac{{81m\sigma_c^2\log^2(1 /\gamma)}}{\bar \epsilon_c^2},~ B_0^v =  \frac{81\sigma_v^2\log^2(1 /\gamma)}{\bar \epsilon_v^2} 
\end{align}
and for $k\ge1,$
\be\label{batch-rm}\begin{aligned}
    B_k^f &= \max \left(\frac{324(L_g^f)^2\|x_k - x_{k-1}\|^2\log^2(1/\gamma)}{\alpha_f\bar \epsilon_f^2},\frac{81\sigma_f^2\log^2(1/\gamma)}{\bar \epsilon_f}\right) , \\
    B_k^c &= \max \left(\frac{{324 m (L_c^f)^2\|x_k - x_{k-1}\|^2\log^2(1/\gamma)}}{\alpha_c\bar \epsilon_c^2},\frac{{81m\sigma_c^2\log^2(1/\gamma)}}{\bar \epsilon_c}\right) , \\
    B_k^v &= \max \left(\frac{324 L_c^2\|x_k - x_{k-1}\|^2\log^2(1/\gamma)}{\alpha_v\bar \epsilon_v^4},\frac{81\sigma_v^2\log^2(1/\gamma)}{\bar \epsilon_v^3}\right) ,
\end{aligned}\ee
where $\alpha_f,\alpha_c,\alpha_v,\bar \epsilon_f,\bar \epsilon_c,\bar \epsilon_v>0$. Then for any $K\ge 1,$ it holds with probability at least $1-6K\gamma$ that for any \(k=0,\ldots, K-1\),
    \be\label{hig-err}\begin{aligned}
        \|\tilde \nabla f_k - \nabla f_k\|^2 &\leq 2\bar \epsilon_f^2+2\alpha_f^{1/2}\bar \epsilon_f^{3/2}+\alpha_f\bar \epsilon_f,\\
        \|\tilde \nabla c_k - \nabla c_k\|^2 &\leq 2\bar \epsilon_c^2+2\alpha_c^{1/2}\bar \epsilon_c^{3/2}+\alpha_c\bar \epsilon_c,\\
        \|\tilde c_k - c_k\|^2 &\leq 2\bar \epsilon_v^4+2\alpha_v^{1/2}\bar \epsilon_v^{7/2}+\alpha_v\bar \epsilon_v^3.
    \end{aligned}\ee
\end{lemma}

\begin{proof}
    For simplicity, we provide the error analysis for the stochastic gradient estimate of the objective, $\|\tilde{\nabla} f_k - \nabla f_k\|^2$, while the results for the other two error terms can be derived similarly. For each $\xi \in \mathcal{B}_k^f$, we define $F_\xi(x) := F(x, \xi)$ for brevity. Define
    \begin{align}
    a_\xi = \nabla F_\xi(x_k) - \nabla F_\xi(x_{k-1}) - \nabla f_k + \nabla f_{k-1}. \notag
    \end{align}
    Then we have  $\mathbb{E}_\xi[a_\xi] = 0$, and the random variables $a_\xi$ are independent and identically distributed, and moreover, 
    \begin{align}
    \|a_\xi\| \leq \|\nabla F_\xi(x_k) - \nabla F_\xi(x_{k-1})\| + \|\nabla f_k - \nabla f_{k-1}\| \leq 2 L_g^f \|x_k - x_{k-1}\|, \notag
    \end{align}
    where the second inequality is  derived from the $L_g^f$-smoothness of both $f$ and $F_\xi$. 
    Applying the Vector Azuma-Hoeffding inequality (Lemma~\ref{lm:vazuma}),  we obtain that with probability at least $1 - \gamma$,  
    \begin{align*}  
        & \left\| \nabla F_{\mathcal{B}_k^f}(x_k) - \nabla F_{\mathcal{B}_k^f}(x_{k-1}) - \nabla f_k + \nabla f_{k-1} \right\| \\  
        &= \frac{1}{B_g^f} \left\| \sum_{\xi \in \mathcal{B}_k^f} \left[ \nabla F_\xi(x_k) - \nabla F_\xi(x_{k-1}) - \nabla f_k + \nabla f_{k-1} \right] \right\|  
        \leq 6 L_g^f \sqrt{\frac{\log(1/\gamma)}{B_g^f}} \|x_k - x_{k-1}\|,
    \end{align*}  
    where $F_{\mathcal{B}_k^f}(x) := \frac{1}{B_g^f} \sum_{\xi \in \mathcal{B}_k^f} F_\xi(x)$.  
    Note that for each $\xi \in \mathcal{B}_k^f$,   
    \( 
    \mathbb{E}[\nabla F_\xi(x_k) - \nabla f_k] = 0, 
    \) 
    and $\|\nabla F_\xi(x_k) - \nabla f_k\| \leq \sigma_f$ by Assumption \ref{ass:stochastic}.  
    Thus, using Lemma~\ref{lm:vazuma} again, it holds with probability at least $1 - \gamma$  that  
    \begin{align}  
        \left\| \nabla F_{\mathcal{B}_k^f}(x_k) - \nabla f_k \right\| = \frac{1}{B_g^f} \left\| \sum_{\xi \in \mathcal{B}_k^f} \left[ \nabla F_\xi(x_k) - \nabla f_k \right] \right\| \leq 3 \sigma_f \sqrt{\frac{\log(1/\gamma)}{B_g^f}}.\notag  
    \end{align}  
    Furthermore, we note that $\tilde \nabla f_k - \nabla f_k = \sum_{t=0}^k (1 - \alpha_f)^t u_{k-t}$, where  
    \begin{align}  
        u_k = \begin{cases}  
         \nabla F_{\mathcal{B}_k^f}(x_k) - \nabla f_k + (1 - \alpha_f) \left( \nabla f_{k-1} - \nabla F_{\mathcal{B}_k^f}(x_{k-1}) \right), & k > 0, \\ 
         \nabla F_{\mathcal{B}_k^f}(x_k) - \nabla f_k, & k = 0.
        \end{cases}\notag  
    \end{align}  
    We have $\mathbb{E}[u_k | x_k] = 0$. 
    For $k = 0$, it holds that with probability at least $1 - \gamma$,  
    \begin{align}  
        \|u_0\| \leq 3 \sigma_f \sqrt{\frac{\log(1/\gamma)}{B_g^0}} \leq \frac{\bar \epsilon_f}{3\sqrt{ \log(1/\gamma)} }\label{gradientVar_1}  
    \end{align}  
    due to the choice of $B_g^0$ in \eqref{rm_b0}.  
    For any $k \ge 1$, conditioned on $x_k$, with probability at least $1 - \gamma$,  we have
    \begin{align}  
        \|u_k\| \leq 3 \alpha_f \sigma_f \sqrt{\frac{\log(1/\gamma)}{B_g^f}} + 6 (1 - \alpha_f) L_g^f \sqrt{\frac{\log(1/\gamma)}{B_g^f}} \|x_k - x_{k-1}\| \leq \frac{(1 - \alpha_f) \alpha_f^{1/2} \bar \epsilon_f + \alpha_f \bar \epsilon_f^{1/2}}{3\sqrt{ \log(1/\gamma)} },\label{gradientVar_0}  
    \end{align}  
    due to the choice of $B_g^f$ in \eqref{batch-rm}.  
    Then by the union bound, the event that \eqref{gradientVar_0} for $k=0$ and \eqref{gradientVar_1} for $1\le k\le K-1$ holds with probability at least $1 - K \gamma$.  
    Conditioned on this event, for any $0\le k\le K-1$ it follows from  Lemma~\ref{lm:vazuma}   that with probability at least $1 - \gamma$,  
    \begin{align}  
        \|\tilde \nabla f_k - \nabla f_k\|^2 &= \left\| \sum_{t=0}^k (1 - \alpha_f)^t u_{k-t} \right\|^2 \notag \\  
        &\leq 9 \log(1/\gamma) \left( \frac{\bar \epsilon_f^2 + 2 \alpha_f^{1/2} \bar \epsilon_f^{3/2} + \alpha_f \bar \epsilon_f}{9 \log(1/\gamma)} + \frac{\bar \epsilon_f^2}{9\log(1/\gamma)} \right) \notag \\  
        &= 2 \bar \epsilon_f^2 + 2 \alpha_f^{1/2} \bar \epsilon_f^{3/2} + \alpha_f \bar \epsilon_f,\label{gradientVar_3}  
    \end{align}  
    where the first inequality leverages $\sum_{t=0}^k (1 - \alpha_f)^{2t} \leq 1/\alpha_f$.  
    Then  by the union bound again, with probability at least $1 - 2K \gamma$ the inequality   \eqref{gradientVar_3} holds for all $k=0,\ldots,  K-1$.   %thus establishing the bound on $\|\tilde{\nabla} f_k - \nabla f_k\|^2$. 
    Similar analysis can be made to $\|\tilde \nabla c_k-c_k\|^2$ and $\|\tilde c_k - c_k\|^2$. Consequently, we derive that with probability at least $1-6K\gamma$ the relations in \eqref{hig-err} hold simultaneously for all $k=0,\ldots,K-1.$%Thus the conclusion can be derived accordingly.
\end{proof}

\end{document}
